\newpage
\chapter{Injection}
The purpose of the following models is to describe accurately the mass flow injection in different conditions.
As for the tank, the type of fluid influences the model to be used: to match the characteristics of a self-pressurising fluid
the Non-Homogeneus Non-Equilibrium (NHNE) model \cite{MassFlowRate2013} was used.
It basically condiders a weighted average of the Homogeneus Equilibrium Model,
with no slip and multi-phase fluid
\begin{equation}
	\dot{m}_{HEM}=A\rho \sqrt{2(h_1-h_2)}
\end{equation}
and the Single Phase Injection
\begin{equation}
	\dot{m}_{SPI}=A \sqrt{2\rho (p_1-p_2)}
\end{equation}
The weight takes into account the pressure drop and the drop respect to the vapor pressure
\begin{equation}
	\chi=\sqrt{\frac{p_1-p_2}{p_V-p_2}}
\end{equation}
and finally the mass flow can be calculated as
\begin{equation}
	\dot{m}_{NHNE}=\frac{\chi}{\chi+1}\dot{m}_{SPI}+\frac{1}{\chi+1}\dot{m}_{HEM}
\end{equation}
This model clearly applies to liquid or multi-phase flow, instead with an ideal gas we write
\begin{equation}
	\dot{m}_{gas}= \frac{p_1A}{\sqrt{RT_1}}f(\frac{p_2}{p_1})
\end{equation}
with
\begin{equation*}
	f(\frac{p_2}{p_1})=\sqrt{\frac{2\gamma}{\gamma -1}((\frac{p_2}{p_1})^\frac{2}{\gamma}-(\frac{p_2}{p_1})^\frac{\gamma +1}{\gamma})}
\end{equation*}
remembering that a chocked exit section is reached when
\begin{equation}
	\frac{p_2}{p_1} \leq (\frac{p_e}{p_1})_{crit}=(\frac{2}{\gamma +1})^\frac{\gamma}{\gamma -1}
\end{equation}
and in that case we write
\begin{equation}
	\dot{m}_{gas}= \frac{p_1A}{\sqrt{RT_1}}\Gamma
\end{equation}
with
\begin{equation*}
	\Gamma=\sqrt{\gamma (\frac{2}{\gamma +1})^\frac{\gamma +1}{\gamma -1}}
\end{equation*}
\section{Functions}
The following functions are in \textit{Injection//PyInjection.py}.
\begin{verbatim}
NHNE_injection(p1, p2, T, cD, fluid)
-> return mdot
\end{verbatim}
\hspace{2em} Calculates mass flux [kg/sm$^2$] using NHNE model.
Returns $0$ if $p_1<p_2$.
 \begin{verbatim}
gas_injection(p1, p2, T, CD, fluid)
-> return mdot
\end{verbatim}
\hspace{2em} Calculates mass flux [kg/sm$^2$] for an ideal gas using CoolProp properties.
Returns $0$ if $p_1<p_2$.
\begin{verbatim}
gas_injection_custom(p1, p2, T, CD, gamma, M, eps=1.0)
-> return mdot
\end{verbatim}
\hspace{2em} Calculates mass flux [kg/sm$^2$] for an ideal gas.
It is meant to be used for nozzle outflow, using burned gas properties, 
so it takes in input also the expansion ratio,modifying the chocked condition to
\begin{equation}
	f(\frac{p_2}{p_1}) \geq \frac{\Gamma}{\varepsilon}
\end{equation}
Returns $0$ if $p_1<p_2$.
\begin{verbatim}
massflow(p1, p2, T, CD, fluid)
-> return mdot
\end{verbatim}
\hspace{2em} Chooses the model to use with a \textit{try/except} structure, because the calculation of vapor pressure using CoolProp will
fail if it is super-critical.\\

All the previous functions take into account the effects of the discharge coefficient.





